\section{Цель
работы}\label{ux446ux435ux43bux44c-ux440ux430ux431ux43eux442ux44b}

Приобретение практических навыков по конфигурированию SMTP-сервера в
части настройки аутентификации

\section{Выполнение лабораторной
работы}\label{ux432ux44bux43fux43eux43bux43dux435ux43dux438ux435-ux43bux430ux431ux43eux440ux430ux442ux43eux440ux43dux43eux439-ux440ux430ux431ux43eux442ux44b}

Во втором терминале сервера запустим вывод логов почты (рис.
{[}-@fig:001{]}).

\begin{figure}
\centering
\pandocbounded{\includegraphics[keepaspectratio,alt={Логи почты /var/log/maillog}]{image/1.png}}
\caption{Логи почты /var/log/maillog}\label{fig:001}
\end{figure}

В первом же терминале сервера авторизуемся под рутом и откроем файл
конфигурации /etc/dovecot/dovecot.conf.

В этом файле допишем в список рабочих протоколов протокол lmtp (рис.
{[}-@fig:003{]}).

\begin{figure}
\centering
\pandocbounded{\includegraphics[keepaspectratio,alt={/etc/dovecot/dovecot.conf}]{image/3.png}}
\caption{/etc/dovecot/dovecot.conf}\label{fig:003}
\end{figure}

Далее, откроем конфигурационный файл /etc/dovecot/conf.d/10-master.conf.
Внутри этого файла пропишем следующее тело для структуры service lmtp
(рис. {[}-@fig:005{]}).

\begin{figure}
\centering
\pandocbounded{\includegraphics[keepaspectratio,alt={service lmtp}]{image/5.png}}
\caption{service lmtp}\label{fig:005}
\end{figure}

Далее, пропишем в postfix сокет, через который будет идти отправка
сообщений. После этого откроем файл /etc/dovecot/conf.d/10-auth.conf
(рис. {[}-@fig:006{]}).

\begin{figure}
\centering
\pandocbounded{\includegraphics[keepaspectratio,alt={Настройка сокета}]{image/6.png}}
\caption{Настройка сокета}\label{fig:006}
\end{figure}

В этом файле зададим значение для поля auth\_username\_format,
отвечающее за формат имени пользователя для аутентификации. В нашем
случае домен не будет указываться (рис. {[}-@fig:007{]}).

\begin{figure}
\centering
\pandocbounded{\includegraphics[keepaspectratio,alt={auth\_username\_format}]{image/7.png}}
\caption{auth\_username\_format}\label{fig:007}
\end{figure}

Перезапустим postfix и dovecot (рис. {[}-@fig:008{]}).

\begin{figure}
\centering
\pandocbounded{\includegraphics[keepaspectratio,alt={Перезапуск postfix и dovecot}]{image/8.png}}
\caption{Перезапуск postfix и dovecot}\label{fig:008}
\end{figure}

Теперь перейдём на виртуальную машину клиента. Попробуем отправить
письмо самому себе (рис. {[}-@fig:009{]}).

\begin{figure}
\centering
\pandocbounded{\includegraphics[keepaspectratio,alt={Отправка письма}]{image/9.png}}
\caption{Отправка письма}\label{fig:009}
\end{figure}

В логах, которые мы открывали на сервере в самом начале выполнения
лабораторной работы, мы видим, что письмо было доставлено в ящик. Об
этом свидетельствует подпись ``saved mail to INBOX''. Кроме того, теперь
в логах пишется, что транспортировка осуществляется через lmtp (passing
\ldots{} to transport=lmtp) (рис. {[}-@fig:010{]}).

\begin{figure}
\centering
\pandocbounded{\includegraphics[keepaspectratio,alt={Логи почты}]{image/1.png}}
\caption{Логи почты}\label{fig:010}
\end{figure}

Откроем на сервере почтовый ящик, чтобы убедится, что письмо успешно
доставлено. Как видим, это действительно так (рис. {[}-@fig:011{]}).

\begin{figure}
\centering
\pandocbounded{\includegraphics[keepaspectratio,alt={Почтовый ящик mail}]{image/11.png}}
\caption{Почтовый ящик mail}\label{fig:011}
\end{figure}

Теперь откроем файл конфигурации по пути
/etc/dovecot/conf.d/10-master.conf.

Приведем содержание тела структуры service auth к следующему виду.
Разберём построчно

\begin{verbatim}
service auth { ... } - Эта строка объявляет начало секции конфигурации для внутренней службы Dovecot, которая называется auth.   

unix_listener /var/spool/postfix/private/auth { ... } - Указывает Dovecot создать "слушателя" на основе UNIX-сокета.   

group = postfix - Устанавливает группу-владельца для файла сокета.   
user = postfix - Устанавливает пользователя-владельца для файла сокета.   
mode = 0660 - Устанавливает права доступа к файлу сокета в восьмеричном формате.   
      
unix_listener auth-userdb { ... } - Создает второй UNIX-сокет.   

mode = 0600 - Устанавливает права доступа для этого внутреннего сокета.   

user = dovecot - Устанавливает пользователя-владельца dovecot.   
\end{verbatim}

(рис. {[}-@fig:013{]}).

\begin{figure}
\centering
\pandocbounded{\includegraphics[keepaspectratio,alt={Структура service auth}]{image/13.png}}
\caption{Структура service auth}\label{fig:013}
\end{figure}

Теперь настроим аутентификацию почты smtp для postfix и укажем, каким
правилам следовать для работы с почтой и её фильтрации. Рассмотрим опции

reject\_unknown\_recipient\_domain - Отклонить письмо, если домен в
адресе получателя не существует.

permit\_mynetworks - Разрешить письмо без дальнейших проверок, если
IP-адрес клиента, отправляющего письмо, находится в списке доверенных
сетей.

reject\_non\_fqdn\_recipient - Отклонить письмо, если адрес получателя
не является полностью определённым доменным именем.

reject\_unauth\_destination - отклоняет письмо, если домен получателя не
является локальным для этого сервера и при этом сессия не
аутентифицирована.

reject\_unverified\_recipient - Отклонить письмо, если Postfix не может
проверить существование получателя.

permit - Если ни одно из предыдущих правил не отклонило и не разрешило
письмо, это правило разрешает его.

(рис. {[}-@fig:014{]}).

\begin{figure}
\centering
\pandocbounded{\includegraphics[keepaspectratio,alt={Настройка почты в postfix}]{image/14.png}}
\caption{Настройка почты в postfix}\label{fig:014}
\end{figure}

Теперь отредактируем файл /etc/postfix/master.cf.

В этом файле внесём изменения так, чтобы smtp поддерживал авторизацию по
sasl (рис. {[}-@fig:016{]}).

\begin{figure}
\centering
\pandocbounded{\includegraphics[keepaspectratio,alt={Включение авторизации}]{image/16.png}}
\caption{Включение авторизации}\label{fig:016}
\end{figure}

Теперь перезапустим postfix и dovecot (рис. {[}-@fig:017{]}).

\begin{figure}
\centering
\pandocbounded{\includegraphics[keepaspectratio,alt={Перезапуск postfix и dovecot}]{image/8.png}}
\caption{Перезапуск postfix и dovecot}\label{fig:017}
\end{figure}

Теперь на клиенте установим telnet (рис. {[}-@fig:018{]}).

\begin{figure}
\centering
\pandocbounded{\includegraphics[keepaspectratio,alt={Установка Telnet}]{image/18.png}}
\caption{Установка Telnet}\label{fig:018}
\end{figure}

Теперь получим ключ авторизации. Этот ключ представляет из себя строку,
содержащую имя пользователя и пароль, и зашифрованную в base64. Теперь
по telnet подключимся к почтовому серверу и проверим соединение. После
этого попробуем с помощью команды auth авторизироваться, в качестве
ключа используя нашу base64 строку. Как видим, авторизация прошла
успешно (Authentification successful) (рис. {[}-@fig:019{]}).

\begin{figure}
\centering
\pandocbounded{\includegraphics[keepaspectratio,alt={Авторизация по telnet}]{image/19.png}}
\caption{Авторизация по telnet}\label{fig:019}
\end{figure}

Теперь настроим сертификаты для postfix, а также уровень security и путь
к базе данных кэша (рис. {[}-@fig:020{]}).

\begin{figure}
\centering
\pandocbounded{\includegraphics[keepaspectratio,alt={Настройка postfix}]{image/20.png}}
\caption{Настройка postfix}\label{fig:020}
\end{figure}

Вновь откроем файл /etc/postfix/master.cf и изменим его следующим
образом (рис. {[}-@fig:021{]}).

\begin{figure}
\centering
\pandocbounded{\includegraphics[keepaspectratio,alt={/etc/postfix/master.cf}]{image/21.png}}
\caption{/etc/postfix/master.cf}\label{fig:021}
\end{figure}

Теперь настроим firewall, разрешив использовать smtp-submission, и
перезапустим postfix (рис. {[}-@fig:022{]}).

\begin{figure}
\centering
\pandocbounded{\includegraphics[keepaspectratio,alt={Настройка firewall}]{image/22.png}}
\caption{Настройка firewall}\label{fig:022}
\end{figure}

Теперь подключимся к серверу через openssl (рис. {[}-@fig:023{]}).

\begin{figure}
\centering
\pandocbounded{\includegraphics[keepaspectratio,alt={openssl}]{image/23.png}}
\caption{openssl}\label{fig:023}
\end{figure}

Теперь сохраним внесённые нами изменения в vagrant (рис.
{[}-@fig:029{]}).

\begin{figure}
\centering
\pandocbounded{\includegraphics[keepaspectratio,alt={vagrant}]{image/29.png}}
\caption{vagrant}\label{fig:029}
\end{figure}

На сервере изменим скрипт mail.sh следующим образом (рис.
{[}-@fig:030{]}).

\begin{figure}
\centering
\pandocbounded{\includegraphics[keepaspectratio,alt={mail.sh для сервера}]{image/30.png}}
\caption{mail.sh для сервера}\label{fig:030}
\end{figure}

\section{Выводы}\label{ux432ux44bux432ux43eux434ux44b}

В результате выполнения лабораторной работы были получены навыки
продвинутой настройки smtp и авторизации

\protect\phantomsection\label{refs}
